\bookchapter{Appendix F: Foldable Phone Purchasing and Usage Special Topic}

\chapterintro{Currently, there are two mainstream categories:}

\booksection{I. How Far Have Foldable Screens Developed?}

: When closed, it's the size of a regular phone; when opened, it's nearly square like a small tablet. Representatives: Huawei Mate X series, Samsung Galaxy Z Fold series, Xiaomi MIX Fold series, vivo X Fold series. Main selling point: “one phone = phone + tablet”.\textbf{Vertical Folding (Clamshell-style)}: When closed, it looks like a compact powder case; when opened, it's a normal-sized phone. Representatives: Huawei Pocket series, Samsung Galaxy Z Flip series, OPPO Find N Flip, Xiaomi MIX Flip. Main selling point: “compact portability + cover screen quick view”.

Key advancements in 2024-2025: Hinge lifespan is generally rated at 500,000 folds (opening/closing 100 times daily lasts 13 years), UTG ultra-thin glass has replaced early plastic films, creases are visible to the naked eye but barely perceptible by touch, prices have dropped from the 10,000+ yuan range down to the 5,000-6,000 yuan range.

\begin{itemize}
  \item \textbf{One-Sentence Verdict}If your current slab phone fully meets your needs and you have no pain point of "screen not being big enough," don't buy a foldable. It's a "nice-to-have" device, not a "must-have" device.
  \item \textbf{II. Who Should Buy a Foldable Phone?}Based on real user feedback, these groups have a high probability of finding foldables "worth it":
\end{itemize}

Key advances in 2024-2026: hinge life generally rated at 500,000 folds (100 folds/day for 13 years), UTG ultra-thin glass replacing early plastic films, creases visible but barely felt by touch, prices dropping from the 10,000 yuan tier to 4,000-5,000 yuan.

\begin{tcbinfo}{: Frequently view PDFs, annotate documents, and read long articles on their phone. When unfolded, the \textasciitilde{}7.5-inch inner screen lets you view A4 PDFs without zooming—a transformative experience.}
Frequent Travelers / Business Users

\end{tcbinfo}

\booksection{II. Who Should Buy a Foldable Phone?}

Creatives Who Love “Flex Mode”

\begin{enumerate}
  \item \textbf{: Fold to 90° and place on a desk—upper screen for framing, lower screen for shutter (selfies/timelapse); upper screen video, lower screen notes; upper screen livestream, lower screen comments—a form factor slab phones simply cannot achieve.}Female Users Who Want Portability + Big Screen
  \item \textbf{: Vertical foldables fit into small bags when closed, unfold to normal phone size, and the cover screen can serve as a makeup mirror or notification viewer.}These People Should NOT Buy One
  \item \textbf{If you only use your phone for short videos, WeChat, and calling—foldables offer zero benefit and are heavier and thicker instead.}If you have small hands or wrist tendonitis—horizontal foldables weigh 230g+ when unfolded, making single-handed use tiring.
  \item \textbf{If you're on a tight budget and plan to keep a phone for 4+ years—foldables have lower resale value than slab phones, and repair costs are high.}III. Buying Guide: Check These 6 Metrics
\end{enumerate}

\begin{tcbwarn}{Don't get dazzled by the word "folding"—focus on these key points when shopping}
\begin{itemize}
  \item Hinge Type
  \item : Waterdrop hinge (shallower crease, longer lifespan, currently mainstream) is better than U-type hinge (deeper crease, early solution). Ask sales staff or check reviews to confirm which type.
  \item Inner Screen Material
\end{itemize}

\end{tcbwarn}

\booksection{III. Buying Guide: Check These 6 Metrics}

Crease Performance

\begin{itemize}
  \item \textbf{: Go to a store and unfold it yourself. The 2025 qualified standard: almost invisible head-on with screen on, visible as a shallow groove under angled light, slight depression when finger runs across but not sharp.}Cover Screen Size
  \item \textbf{: For horizontal foldables, the larger the outer screen the better (ideally 6.4"+) so it works as a normal phone when closed; for vertical foldables, larger is more practical (4"+ allows replying to messages and viewing navigation).}Weight \& Thickness
  \item \textbf{: Horizontal foldables at 230-260g unfolded is normal; under 230g is excellent. Closed thickness of 12-14mm; under 12mm is thin. Vertical foldables must fit comfortably in a shirt pocket when closed to qualify.}: Go to a physical store and unfold it yourself. The 2026 acceptable standard is: almost invisible when screen is on, a shallow groove visible under angled light, slight depression when fingers pass over but not sharp.
  \item \textbf{: Focus on inner screen crack repair cost. Some brands offer free inner screen replacement once in the first year; others charge 2,000-4,000 yuan out-of-pocket. This is the biggest hidden cost of owning a foldable.}In-Store Experience Checklist
  \item \textbf{Open and close it 10 times yourself—feel if the hinge damping is smooth and if there are any unusual sounds.}Unfold it and look at the crease under angled light, then look at it head-on with the screen on.
  \item \textbf{Hold it single-handed in unfolded state for 30 seconds—feel if the weight is acceptable.}Try Flex Mode: fold to 90°, place on desk, try taking photos and watching videos.
\end{itemize}

\begin{tcbpractice}{Ask clearly about the standalone inner screen repair price and warranty policy.}
\begin{enumerate}
  \item IV. Unique Foldable Usage Tips
  \item After buying one, here are experiences you can't get with a slab phone:
  \item 1. Split-Screen \& Multitasking
  \item When a horizontal foldable is unfolded, most systems support top/bottom or left/right split-screen. Typical scenarios: upper screen for online course video, lower screen for notes; left side WeChat, right side documents; upper screen map navigation, lower screen music. Much more practical than slab phone split-screens because the screen is large enough that neither half feels cramped.
  \item 2. Flex Mode (90° Fold)
\end{enumerate}

\end{tcbpractice}

\booksection{IV. Unique Foldable Usage Tips}

Selfies / Group Photos

\booksubsection{: Place on desk folded to 90°, upper screen for framing, lower screen for shutter—no need to ask someone else to take the photo.}

Timelapse / Night Shots

\booksubsection{: Same stable placement on desk—saves you from needing a tripod.}

Video Calls

\begin{itemize}
  \item \textbf{: Frees both hands; the other person sees you from a more natural angle too.}Watching Videos
  \item \textbf{: Use as a mini TV on a desk—especially practical on airplane/high-speed rail tray tables.}Cover Screen Preview + Inner Screen Shooting
  \item \textbf{: When using the rear camera to shoot yourself, the cover screen displays the live preview for more accurate composition.}3. Cover Screen Quick View
  \item \textbf{When a vertical foldable is closed, the cover screen can: show time/notifications/weather, reply to WeChat (voice-to-text or quick replies), tap transit cards, view navigation, act as a makeup mirror (using front camera). No need to unfold every time—saves power and effort.}4. Desktop Mode
  \item \textbf{Some horizontal foldables (Huawei, Samsung) can switch to "Desktop Mode" when unfolded. After connecting an external display, keyboard, and mouse, it works like a basic PC—adequate for temporary document processing and email replies. Not a primary productivity tool, but great for emergencies.}V. Maintenance Tips: Help Your Foldable Last Longer
\end{itemize}

\booksubsection{The most delicate parts of a foldable are the inner screen and hinge. These habits can significantly extend its lifespan:}

Don't Force Open/Close in Low Temperatures

\booksubsection{: Below 0°C, UTG glass becomes brittle. Warm it in your pocket for a few minutes before opening/closing. This is something written in every official manual but nobody reads.}

Avoid Nails/Hard Objects Touching the Inner Screen

\booksection{V. Maintenance Tips: Help Your Foldable Last Longer}

Check for Foreign Objects Before Closing

\begin{itemize}
  \item \textbf{: A single grain of sand caught inside when closing could directly crush the inner screen. Glance over it before closing.}Don't Apply Regular Tempered Glass Film
  \item \textbf{: The inner screen can only use official specialized film or UTG protective film. Regular tempered glass film affects folding and touch functionality. Recommend going to official after-sales for screen protector application.}Keep Dust \& Liquids Away from the Hinge
  \item \textbf{: Although most claim IPX8 water resistance, dust resistance ratings are generally low. Be careful at beaches, construction sites, etc.}Apply Even Force When Opening/Closing
  \item \textbf{: Don't violently flick it open with one hand. Opening smoothly with both hands from the center is most stable.}High-Risk Scenarios for Inner Screen Damage
  \item \textbf{Catching headphone/charging cable when closing—directly cracks the inner screen.}Mixing with keys in a bag—outer screen fine, but keys press creases into the inner screen.
  \item \textbf{Immediately opening/closing when moving from cold outdoors to warm indoors in winter—temperature difference causes glass stress changes.}After a year of use, feeling "the crease got deeper" and trying to peel off the factory film—peeling the film will damage the inner screen along with it.
\end{itemize}

\textbf{VI. Foldable “Pitfalls” \& Realistic Expectations}\par

\begin{itemize}
  \item Before buying, keep these things in mind:
  \item The Crease Will Always Be There
  \item : 2025 technology cannot achieve "completely invisible"—only "almost invisible head-on with screen on". If you're a perfectionist who would stare at the crease uncomfortably every day, don't buy.
  \item Heavier \& Thicker
\end{itemize}

\booksection{VI. Foldable “Pitfalls” \& Realistic Expectations}

Battery Life Is Generally Average

\begin{itemize}
  \item \textbf{: Because space must accommodate the hinge mechanism, battery capacity is compressed. At the same price point, battery life typically falls short of slab phones.}: 2026 technology can't achieve "completely crease-free" — only "almost invisible from the front when screen is on." If you're a perfectionist who'll be bothered by the crease daily, don't buy.
  \item \textbf{: After one year, the second-hand price is usually 50-60\% of the new phone price—drops faster than slab phones.}Expensive Repairs
  \item \textbf{: Inner screen assembly repair costs 2,000-4,000 yuan; hinge repair costs 500-1,500 yuan. Check warranty policies clearly before buying.}Inconsistent App Adaptation
  \item \textbf{: Mainstream apps have adapted to large screens, but niche apps may just be "stretched/enlarged" with mediocre experience.}A foldable is not a "better phone"—it's "a phone in a different form factor." Buy it because you need its unique form, not because it's "more advanced."
  \item \textbf{—— On the Philosophy of Choosing a Foldable Phone}VII. How to Choose Among Mainstream Brands (2025)
  \item \textbf{Not listing specific models (they'll become outdated)—just brand personalities to help you quickly orient yourself:}Huawei Mate X Series
\end{itemize}

: The benchmark for horizontal folding. Strong HarmonyOS ecosystem synergy. Cover screen ratio close to a normal phone—works as a primary device when closed. High price but relatively best resale value.

Samsung Galaxy Z Fold/Flip

\booksection{VII. How to Choose Among Mainstream Brands (2026)}

Xiaomi MIX Fold/Flip

\begin{itemize}
  \item \textbf{: Value-for-money approach. Packed with specs at relatively affordable prices. Suitable for users who want to try foldables without spending too much.}vivo X Fold Series
  \item \textbf{: Among horizontal foldables, the cover-to-inner screen ratio comes closest to a normal phone—great experience using it as a slab phone when closed.}OPPO Find N Series
  \item \textbf{: Among horizontal foldables, the size leans compact—single-hand holdable, suitable for users who find large foldables too heavy.}Huawei Pocket Series
  \item \textbf{: Among vertical foldables, the largest and most practical cover screen. Makeup mirror, UV detection, and other women-oriented features are well-implemented.}Buying advice: If budget allows and you want ecosystem integration, choose Huawei Mate X; for the most mature hinge, choose Samsung Z Fold; for best value, choose Xiaomi/vivo; for vertical folding, look at Huawei Pocket and Samsung Z Flip.
  \item \textbf{Final Reminder}Foldable screen technology iterates every year. This appendix was written in 2025. Before buying, be sure to check the latest reviews and user feedback, focusing especially on: actual hinge lifespan testing, crease performance, and inner screen repair policy. Hands-on experience in-store always matters more than reading reviews.
  \item \textbf{Last read at: ...}Continue Reading
\end{itemize}

Dismiss

\begin{tcbinfo}{Final Reminder}
Foldable screen technology iterates every year. This appendix was written in 2026. Before buying, be sure to check the latest reviews and user feedback, focusing on: hinge life tests, crease performance, and inner screen repair policies. Touching it in person at a store is always more important than reading reviews.

\end{tcbinfo}
